\chapter{Tecnologias}

Este capítulo apresenta as principais tecnologias empregadas no desenvolvimento
do aplicativo, organizadas de acordo com a camada da arquitetura em que
atuam: front-end, back-end e infraestrutura/persistência de dados.

\section{Front-end}

\subsection{React Native}

React Native é um \textit{framework} de código aberto, mantido pela Meta, que
permite o desenvolvimento de aplicativos móveis multiplataforma a partir de
uma única base de código escrita em JavaScript/TypeScript, utilizando
componentes que são traduzidos para elementos nativos tanto no iOS quanto no
Android \citep{reactnative:2026}.

A escolha do React Native se justifica pelo escopo reduzido da equipe de
desenvolvimento (um único desenvolvedor) e facilidade de conhecimento e
manutenção, sem trazer, para este projeto, benefícios que justifiquem esse
custo adicional. Além disso, o ecossistema maduro do React Native oferece
ampla disponibilidade de bibliotecas para funcionalidades já previstas no
aplicativo, como notificações locais e gráficos.

\subsection{Zustand}

Zustand é uma biblioteca leve de gerenciamento de estado para aplicações
React e React Native, que dispensa a estrutura verbosa de \textit{providers}
e \textit{reducers} tipicamente associada a bibliotecas como Redux
\citep{zustand:2026}.

Ela foi escolhida para gerenciar o estado efêmero da aplicação (por exemplo,
dados de formulários em preenchimento ou estados de navegação temporários)
por sua API minimalista e baixo \textit{overhead} de configuração, adequada a um
projeto de escopo individual em que a complexidade adicional de soluções
mais robustas de gerenciamento de estado não se justifica.

\section{Back-end}

\subsection{Express.js}

Express.js é um \textit{framework} minimalista para aplicações web em
Node.js, amplamente utilizado para a construção de APIs \textit{REST},
oferecendo uma camada fina de abstração sobre o roteamento de requisições HTTP e
\textit{middlewares} \citep{express:2026}.

Sua adoção se deve à simplicidade e à flexibilidade não opinativa do
\textit{framework}, que permite estruturar a API de acordo com as
necessidades específicas do projeto, sem impor convenções rígidas de
arquitetura, além de sua ampla documentação e maturidade no ecossistema
Node.js.

\subsection{Zod}

Zod é uma biblioteca de validação e definição de esquemas de dados para
TypeScript, que permite declarar formatos esperados de entrada (como corpos
de requisição) e inferir automaticamente os tipos estáticos correspondentes
\citep{zod:2026}.

Zod foi escolhido para validar os dados recebidos pela API antes de seu
processamento, garantindo que registros de humor, sono e demais entradas do
usuário estejam em conformidade com o formato esperado. A inferência
automática de tipos a partir dos esquemas evita a duplicação de definições de
tipo entre a camada de validação e o restante do código TypeScript do
back-end.

\subsection{Valkey e BullMQ}

Valkey é um banco de dados em memória, do tipo chave-valor, criado como
\textit{fork} de código aberto do Redis, mantido sob a governança da Linux
Foundation \citep{valkey:2026}. BullMQ é uma biblioteca para gerenciamento de
filas de tarefas em Node.js, construída sobre bancos de dados compatíveis com
o protocolo Redis \citep{bullmq:2026}.

A combinação das duas ferramentas foi adotada para o processamento paralelo
e assíncrono de tarefas que não precisam bloquear a resposta imediata da API,
como o disparo de notificações programadas e o processamento de indicadores
a partir do histórico de registros do usuário. A escolha do Valkey em
específico, e não do Redis diretamente, se deve à mudança de licenciamento do
Redis em 2024, que passou a restringir certos usos comerciais; o Valkey
mantém compatibilidade total de protocolo, permanecendo sob licença
permissiva de código aberto.

\section{Infraestrutura e Persistência de Dados}

\subsection{PostgreSQL}

PostgreSQL é um sistema gerenciador de banco de dados relacional, de código
aberto, reconhecido por sua conformidade com o padrão SQL, robustez e suporte
a tipos de dados avançados \citep{postgresql:2026}.

Sua escolha se justifica pela natureza estruturada e relacional dos dados do
domínio da aplicação (registros de humor associados a usuários, horários,
gatilhos e categorias), que se beneficia das garantias de integridade
referencial e das transações ACID oferecidas por um banco relacional
maduro, em detrimento de soluções não relacionais.

\subsection{Prisma}

Prisma é um \textit{Object-Relational Mapper} (ORM) para o ecossistema
Node.js/TypeScript, que gera automaticamente um cliente de acesso ao banco de
dados tipado a partir de um esquema declarativo, além de oferecer um sistema
de migrações versionadas \citep{prisma:2026}.

O Prisma foi escolhido por permitir a manutenção do esquema do banco de dados
de forma versionada e sincronizada com o código da aplicação, reduzindo a
chance de divergência entre a estrutura do banco e as entidades manipuladas
pelo back-end, além de prover checagem de tipos em tempo de compilação para
as consultas realizadas.

\section{Implantação e Automação}

\subsection{Docker e Docker Compose}

Docker é uma plataforma de conteinerização que permite empacotar uma
aplicação e suas dependências em uma unidade isolada e portátil, capaz de ser
executada de forma consistente em diferentes ambientes \citep{docker:2026}. O
Docker Compose é uma ferramenta complementar que permite definir e orquestrar
múltiplos contêineres -- e as relações entre eles -- por meio de um único
arquivo de configuração declarativo \citep{dockercompose:2026}.


No projeto, o Docker Compose é utilizado tanto no ambiente de desenvolvimento
quanto no ambiente de produção para orquestrar os serviços de banco de dados
(PostgreSQL), armazenamento em memória (Valkey) e da própria API, garantindo
que as versões e configurações desses serviços sejam idênticas entre as
máquinas de desenvolvimento e o servidor de produção.

O uso de \textit{healthchecks} em conjunto com a diretiva
\texttt{depends\_on} e a condição \texttt{service\_healthy} assegura
que a API só seja iniciada após o banco de dados e o Valkey estarem de
fato prontos para aceitar conexões, evitando falhas de inicialização
por condição de corrida (\textit{race condition}), como ilustrado no
Código~\ref{lst:healthcheck-compose}.

\begin{lstlisting}[caption={Configuração dUndefined referencee \textit{healthcheck} do Valkey e de dependência condicional da API no \texttt{docker-compose.yml}}, label={lst:healthcheck-compose}]
valkey:
  image: valkey/valkey:9.0-alpine
  healthcheck:
    test: ["CMD", "valkey-cli", "ping"]
    interval: 10s
    timeout: 5s
    retries: 5
    start_period: 5s

api:
  build: ./api
  depends_on:
    valkey:
      condition: service_healthy
    db:
      condition: service_healthy
\end{lstlisting}

\sourcecode{Elaborado pelo autor (2026)}

A adoção do Docker se justifica, no contexto de uma aplicação de saúde
mental de código aberto, pelo objetivo de garantir a reprodutibilidade do
ambiente de execução: qualquer pessoa que clone o repositório --- seja para
contribuir com o projeto, auditar seu funcionamento ou realizar sua própria
implantação --- consegue reproduzir o ambiente completo sem depender de
configurações manuais específicas do sistema operacional hospedeiro, o que
reforça o caráter \textit{self-service} do projeto.

\subsection{GitHub Actions}

GitHub Actions é uma plataforma de integração e entrega contínuas (CI/CD)
nativa do GitHub, que permite a definição de fluxos de trabalho automatizados
--- disparados por eventos como \textit{push} ou \textit{pull request} ---
descritos de forma declarativa em arquivos YAML \citep{githubactions:2026}.

No projeto, o fluxo de trabalho de CI/CD é dividido em três etapas
sequenciais e dependentes entre si. A primeira etapa (\texttt{test}) sobe um
serviço efêmero de PostgreSQL e executa a suíte de testes automatizados da
API a cada \textit{push} ou \textit{pull request} que modifique arquivos do
diretório da API.

A segunda etapa (\texttt{build-and-push}), condicionada ao
sucesso da primeira e restrita a eventos de \textit{push}, constrói a imagem
Docker de produção da API e a publica no GitHub Container Registry (GHCR),
com marcação automática pelo SHA do \textit{commit} e, condicionalmente, pela
tag \texttt{latest} quando o \textit{push} ocorre na \textit{branch}
\texttt{master}.

Por último, a terceira etapa (\texttt{deploy-production}), restrita à
\textit{branch} \texttt{master}, conecta-se ao servidor de produção via SSH e
realiza a atualização do serviço da API por meio da atualização da tag da
imagem no arquivo \texttt{.env} e da recriação do contêiner correspondente
via Docker Compose.

Esse processo elimina a necessidade de intervenção manual para a maior
parte do ciclo de implantação, restringindo a ação humana direta sobre o
servidor de produção a situações excepcionais, o que reduz a superfície de
erro humano e aumenta a confiabilidade do processo de entrega de novas
versões da aplicação.

\subsection{Ansible}

Ansible é uma ferramenta de automação de infraestrutura que permite
descrever, de forma declarativa e idempotente, o estado desejado de um ou
mais servidores, dispensando a necessidade de um agente instalado na máquina
gerenciada, uma vez que se comunica com ela via SSH \citep{ansible:2026}.

No escopo deste projeto, o papel do Ansible é ortogonal ao do Docker Compose
e ao do GitHub Actions: enquanto estes automatizam, respectivamente, a
orquestração dos serviços da aplicação e o ciclo de integração e entrega
contínuas, o Ansible é responsável exclusivamente pelo provisionamento
inicial do servidor de produção, executado uma única vez --- ou sempre que um
novo servidor precisar ser provisionado. Esse provisionamento consiste em:
(1) criar o usuário dedicado à aplicação; (2) adicionar a chave pública SSH
utilizada pelo GitHub Actions para se autenticar no servidor durante a etapa
de implantação; (3) criar os diretórios da aplicação com as permissões
adequadas; e (4) instalar o Docker no servidor. A cópia do arquivo de
variáveis de ambiente de produção (\texttt{.env.production}), por conter
segredos e credenciais sensíveis, é deliberadamente mantida como uma etapa
manual, não automatizada, e não é versionada no repositório.

A separação entre o provisionamento inicial (\textit{Ansible}), a orquestração de
execução (\textit{Docker Compose}) e a automação do ciclo de
entrega (\textit{GitHub Actions}) reflete uma decisão arquitetural:
cada ferramenta atua em uma camada distinta e bem delimitada do processo
de implantação, o que favorece a reprodutibilidade de todo o processo,
do zero, por qualquer pessoa que deseje hospedar sua própria instância da aplicação.

Essa preocupação é particularmente relevante no contexto de um
aplicativo de saúde mental de código aberto: ao reduzir a dependência
de serviços proprietários de implantação (\textit{walled gardens}) e
ao documentar e automatizar o processo de provisionamento, o projeto
favorece um modelo mais alinhado aos princípios de privacidade e de
autonomia sobre os dados dos usuários, permitindo que instituições ou
indivíduos preocupados com a confidencialidade de dados sensíveis de
saúde mental hospedem e controlem sua própria instância da aplicação.
